\section{Relation to Prior Work}
\label{sec:related}

Cooley and Tukey established the recursive decomposition now associated with
complex-radix FFTs \cite{cooley1965}.  Bruun instead retained real conjugate
factors and interpreted the recursion through real $z$-transform filters
\cite{bruun1978}.  The construction here retains Bruun's half-angle factor tree
and supplies explicit inverse interpolation, normalized quotient coordinates,
and boundary maps.

Chen, Sorensen, and Jones specialized Bruun's structure to real-symmetric data
\cite{chen1992}.  The normalized quotient map here instead treats arbitrary
real input and derives its inverse and metric from the complete real CRT.

Chinese-remainder formulations of Fourier algorithms appear in the work of
Good, Winograd, and the algebraic signal-processing literature
\cite{good1958,winograd1978,puschel2008,voronenko2009}.  In that language,
Cooley-Tukey and Bruun are instances of a common polynomial factorization
scheme expressed in different intermediate bases \cite{voronenko2009}.  The
present normalization makes this basis choice explicit at every real quadratic
leaf and follows its metric consequence through the complete product.

Real-input FFT algorithms commonly exploit conjugate symmetry to reduce
arithmetic and storage \cite{sorensen1987,duhamel1990}.  The packed transform
\eqref{eq:packed-dft} uses the same spectral degrees of freedom, but derives
them directly from real quotient algebras.  The ODFT is related to shifted and
half-bin Fourier grids; its contribution here is closure of the same normalized
two-plane arithmetic on the factors of $z^N+1$.

Stockham schedules and DIT/DIF variants reorder intermediate Fourier states
\cite{stockham1966,duhamel1990}.  The phase-packet result states the common
principle as an exact conjugation of each stage, with the DIP schedule selected
by a lattice shear.  The positive-cone result is independent of a particular
schedule: it lifts any signed stage factorization and preserves the complete
product under rail-difference projection.

\begin{table*}[t]
\caption{Algebraic roles of the transforms developed here}
\label{tab:relation}
\centering
\scriptsize
\begin{tabularx}{\textwidth}{@{}lXXXXX@{}}
\toprule
Object & Root polynomial & Interior state & Boundary order & Exact inverse & Metric\\
\midrule
Normalized Bruun DIF & $z^N-1$ & real CRT residues & native leaves & local CRT ascent & $B^{\trans}B=NI$\\
Bruun DIT & $z^N-1$ & spectral combs & natural spectrum & reversed cells & $B^{\trans}B=NI$\\
DIP & $z^N-1$ & sheared Zak packets & natural spectrum & reversed shear & $B^{\trans}B=NI$\\
ODFT & $z^N+1$ & real half-bin pairs & conjugate pairs & local pair inverse & $O^{\trans}O=NI$\\
Cone lift & either & nonnegative rails & projected boundary & signed inverse after projection & gauge-dependent\\
\bottomrule
\end{tabularx}
\end{table*}
